\documentclass[11pt]{article}
\usepackage{amsmath}
\usepackage{booktabs}
\usepackage[margin=1in]{geometry}

\title{Tau Ceti's Debris Disc: Archival Raw-Data Reduction Phase}
\author{Reproducible research repository}
\date{\today}

\begin{document}
\maketitle

\begin{abstract}
This document records a conditional revision of the accepted Tau Ceti debris-disc
constraint synthesis. The archival phase preserved and checksummed the official
Herschel and ALMA deliveries, performed an executable map-level reanalysis of
the official Herschel products, and executed the delivered ALMA calibration
scripts through a documented modular CASA compatibility layer. Calibrated 12-m
and ACA measurement sets were used for a diagnostic visibility-domain fit and a
joint map/visibility posterior. No new telescope observations were obtained.
Full HIPE telemetry reprocessing and original CASA 4.2.2 binary validation were
not available, so the accepted baseline remains the authoritative headline
inference.
\end{abstract}

\section{Status relative to the baseline}

The accepted baseline gives
\[
f_\tau=8.36\times10^{-6},\qquad
q_\tau=0.927,\qquad
z_\tau=1.46,
\]
with validated collisional ratios \(f/f_{\max}=0.19\)--\(0.23\).
These values are frozen in
\texttt{results/baseline/accepted\_baseline\_manifest.json}.

\section{Archival data and honest reduction level}

The exact Herschel OBSIDs are 1342199389, 1342213575 and 1342213576.
The ALMA project is 2013.1.00588.S with member OUS entries
\texttt{uid://A001/X121/X3b7} and \texttt{uid://A001/X121/X3b9}.
The official raw and product bundles are preserved under
\texttt{data/raw/} and \texttt{data/archive\_products/}.

HIPE is unavailable locally, so the Herschel result is explicitly a map-level
reanalysis of official Level 2/2.5 products rather than a raw telemetry
re-reduction. The delivered CASA 4.2.2/Python-2 scripts were translated in
memory and executed with modular CASA 6.7.5. This produced calibrated
measurement sets and a visibility-domain diagnostic, but is not a claim of
bit-for-bit original-runtime equivalence.

\section{Herschel map-level result}

The reference fit uses deduplicated PACS 70 and 160 micrometre JSCANAM products
and the SPIRE 250 micrometre extended-source map. It fits a PSF-convolved
axisymmetric belt, fixed photospheres, a candidate background source, a
background plane and residual astrometric offsets. The reference geometry is
approximately \(R_{\rm in}=16.6\) au, \(R_{\rm out}=52.3\) au, inclination
\(21^\circ\), and position angle \(93^\circ\). Map-maker and PSF variants shift
these values materially, and injection recovery does not robustly recover the
inner edge. Consequently these numbers are diagnostic map-level estimates,
not replacement posterior constraints.

\section{ALMA and joint inference}

The archive contains the official calibration scripts, raw ASDM deliveries,
QA products and restored images. The 12-m and both ACA execution blocks were
calibrated; 2.27 million 12-m and 0.51 million ACA target visibilities were
extracted and fitted in the uv domain. A diagnostic joint map/visibility
posterior gives \(f_\tau=2.85\times10^{-6}\) (68\% interval
\(1.85\)--\(4.22\times10^{-6}\)), with ring and broad-belt median
\(f/f_{\max}=0.67\) and \(0.74\), respectively. These values are reported in
\texttt{report/joint\_map\_visibility\_posterior.md}; they are not promoted over
the accepted baseline because the Herschel reduction is map-level and CASA is
not the original archive binary.

\section{New-observation proposal}

The accompanying proposal package recommends a compact ALMA Band 6
12-m+ACA continuum program with maximum recoverable scale above \(30''\),
approximately \(0.8''\)--\(1.0''\) 12-m resolution, and a target rms of
10--15 microJy per beam. It is a design study only; no observing time has
been requested or awarded.

\end{document}
